% ---- The problem ----

In multi-quadrotor cooperative cable-suspended payload transport, $N$ underactuated vehicles are connected to a shared rigid payload by flexible cables, forming a system on $\mathcal{Q} = \SEthree \times (\SEthree \times \Sph^2)^N$ with $6 + 8N$ degrees of freedom. Geometric control methods~\cite{lee2010geometric,sreenath2013geometric,lee2018geometric,wu2014geometric} provide almost-global exponential stability on $\SOthree$ and aggressive trajectory tracking, bringing multi-vehicle cable transport close to deployment readiness~\cite{estevez2024review,sun2025agile}. Yet each cable remains a single mechanical point of failure. When a cable snaps, the payload accelerates downward (${\approx}3.3$~m/s$^2$), the freed agent accelerates upward (${\approx}6.5$~m/s$^2$), and each remaining agent experiences an instantaneous load-share jump---a $50\%$ increase for $N{=}3$. These effects propagate from sub-$10$~ms (rigid-body force transmission) to ${\sim}200$~ms (cable wave propagation), producing a multi-channel fault signature that surviving agents must absorb without centralized coordination or inter-agent communication.

This challenge is sharpened by the architecture's defining property. The baseline GPAC (Geometric Position and Attitude Control) hierarchy enforces an \emph{N-agnostic constraint}: each agent knows only its own mass and inertia, local sensor measurements, cable tension and direction, and a shared reference trajectory. It does not know the team size, payload mass, or any other agent's state. This removes single points of failure and communication dependency, but also rules out model-based fault detection such as multiple-model adaptive estimation, which requires $N{+}1$ hypothesis filters. The tension-channel signal-to-noise ratio for a remote cable break is only about~$1.3$, indistinguishable from wind perturbations. Cable-fault tolerance in this setting must therefore work without detection, isolation, or recovery.

% ---- The gap ----

Despite mature geometric transport theory and extensive UAV fault-tolerance work on actuator and sensor failures~\cite{sun2020incremental,zhang2024fault,adaika2025fdd}, no published work addresses cable-fault tolerance in cooperative cable-suspended aerial transport. The closest result, Liang~et~al.~\cite{liang2021fault}, assumes known failure identity and centralized reconfiguration, both incompatible with N-agnostic architectures. Zeng~et~al.~\cite{zeng2025decentralized} and Gao~et~al.~\cite{gao2024decentralized} observed implicit resilience to agent loss in learning-based transport, but did not formalize the mechanism or address cable-suspended underactuated systems.

% ---- The key observation ----

The research objective emerged from an empirical observation. The predecessor system Tether\_Lift~\cite{hajieghrary2026tetherlift} established that the GPAC hierarchy can coordinate cooperative lift using only local measurements under the N-agnostic constraint. The present paper (Tether\_Grace) asks whether the same architectural property also yields resilience to abrupt cable-topology change.

The answer lies in PID feedback with gravity feedforward. Under cable-snap faults with no explicit fault handling, the GPAC baseline maintained stable post-fault tracking with payload RMSE of $44$--$49$~cm across $N \in \{3, 5, 7\}$, decreasing with team size despite cascading faults. The mechanism is the N-agnostic PID structure: each agent uses gravity feedforward from its own mass, PID feedback from its position error, and anti-swing damping from local cable angular velocity. None of these terms depends on $N$, $m_L$, or any other agent's state. When a cable snaps, the surviving agents' tracking errors change through the physics, and the PID integral term absorbs the higher load share without adaptation or re-convergence delay. Ablation confirms that disabling cable-tension feedforward changes RMSE by less than $1\%$ (Section~\ref{sec:tension_ablation}).

This shifts the research problem from \emph{how to detect and recover from cable faults} to \emph{whether the emergent fault tolerance can be formalized and quantified}: under what conditions does an N-agnostic controller with local cable-tension feedforward produce bounded tracking degradation across cable topologies? Figure~\ref{fig:paper_scope} summarizes this narrowing from the broader Tether\_Lift baseline to the present Tether\_Grace fault-tolerance question and its current evidence boundary.

% ---- The approach ----

This paper addresses that question by identifying the structural property that makes the architecture topology-invariant, deriving operational bounds, and validating them with comprehensive simulation evidence. The replayed full-Drake path uses the GPAC hierarchy with PID-plus-feedforward control. A post-fault geometry redistribution path is also described. Adaptive extensions ($\mathcal{L}_1$, concurrent learning) are implemented in the codebase but remain future integration work.

Two structural properties make the PID-plus-gravity-feedforward design attractive. First, PID feedback with gravity feedforward provides the core fault-absorption capability, as confirmed by ablation (Section~\ref{sec:tension_ablation}). Second, all controller terms depend only on the agent's own state, so no agent's control law is invalidated by a topology change it cannot observe.

These properties motivate a \emph{topology-invariance design hypothesis}: if the post-fault disturbance remains within a fault-inclusive bound and the local thrust margin is adequate, then the tracking bound is governed by controller bandwidth and disturbance size rather than cable identity. Full-Drake multibody results support this: under matched single-fault conditions (identical fault, varying $N$), payload RMSE decreases from $48.8$~cm at $N{=}3$ to $38.4$~cm at $N{=}7$. A thrust-saturation boundary sweep identifies a sharp phase transition at $f_{\max} \approx 45$~N that coincides with the theoretical post-fault hover minimum, with the default $150$~N providing $3\times$ headroom.

The contribution of this paper is threefold:

\begin{enumerate}[label=C\arabic*]
  \item \emph{Topology-invariant error structure and mechanism identification:} A structural analysis (Proposition~1) showing that the N-agnostic PID-plus-feedforward control law produces per-agent error dynamics that are identical across cable topologies: cable snaps change the disturbance input, not the system matrices. This motivates a topology-invariance hypothesis that is validated by full-Drake simulation. A finite-time transient prediction (${\sim}18$~cm predicted, observed ${\sim}107$~cm; the ${\sim}5.8\times$ ratio is topology-invariant across scenarios) and a closed-form cost ratio $\rho_{\mathrm{FT}} = N/(N{-}k)$ provide actuator-sizing criteria. Ablation confirms PID feedback with gravity feedforward as the primary mechanism (cable-tension feedforward $< 1\%$ RMSE change).
  \item \emph{Comprehensive multibody evidence:} Full-Drake simulation across $N \in \{3, 5, 7\}$ with $1$--$3$ cascading cable-snap faults under ideal and sensor-realistic (ESKF + Dryden wind) conditions. Monte Carlo validation ($30$ runs per scenario) with randomized fault timing, wind, cable lengths, and initial conditions. Matched single-fault experiments ($48.8 \to 38.4$~cm), fault-count isolation, trajectory generalization, and baseline comparison provide deconfounding evidence.
  \item \emph{Quantitative design framework:} A fault-tolerance cost ratio $\rho_{\mathrm{FT}}$ validated against empirical nominal-vs-fault RMSE ratio, a thrust-margin screening methodology validated by the $f_{\max}$ phase transition at the theoretical minimum, and a tilt-corrected capacity analysis for actuator sizing.
\end{enumerate}

The rest of the paper is organized as follows: Section~\ref{Section_II:Related_Works} reviews related work. Section~\ref{Section_III:Problem_Statement} defines the system model, the cable-snap fault event, and the formal problem. Section~\ref{Section_IV:Fault-Tolerant_Control_Architecture} presents the fault-tolerant architecture and its integration boundaries. Section~\ref{Section_V:Experimental_Methodology} defines the evidence methodology. Section~\ref{Section_VI:Simulation_Results} reports simulation results. Section~\ref{Section_VII:Discussion} interprets findings, identifies gaps, and concludes.

\begin{figure}[t]
\centering
\begin{tikzpicture}[
  node distance=6mm and 8mm,
  block/.style={draw, rounded corners, thick, align=center, minimum width=3.1cm, minimum height=1.0cm, fill=blue!8},
  focus/.style={draw, rounded corners, thick, align=center, minimum width=3.1cm, minimum height=1.0cm, fill=orange!12},
  verdict/.style={draw, rounded corners, thick, align=center, minimum width=3.1cm, minimum height=1.0cm, fill=green!10},
  arrow/.style={-{Latex[length=2.5mm]}, thick}
]
\node[block] (nominal) {Nominal GPAC\\tracking baseline};
\node[block, right=of nominal] (faults) {Cable snap and\\degradation faults};
\node[focus, below=of nominal] (reduced) {Reduced-order\\breadth studies};
\node[focus, below=of faults] (drake) {Matched Drake\\fidelity replay};
\node[verdict, below=10mm of $(reduced)!0.5!(drake)$] (unknowns) {Replayable U1--U6\\status verdicts};

\draw[arrow] (nominal) -- (faults);
\draw[arrow] (nominal) -- (reduced);
\draw[arrow] (faults) -- (drake);
\draw[arrow] (reduced) -- (unknowns);
\draw[arrow] (drake) -- (unknowns);
\end{tikzpicture}
\caption{Overview of the paper scope. The study begins from the validated GPAC baseline (Tether\_Lift), introduces the cable-snap fault problem, documents the Tether\_Grace fault-tolerance extension, and evaluates the present evidence.}
\label{fig:paper_scope}
\vspace{-5pt}
\end{figure}
